\section{Benchmark setup and data}
\label{sec:setup}

\paragraph{Tasks and simulator.} Both tasks run in robosuite \rf{ver.robosuite} (commit
\rf{ver.robosuite.commit}) on MuJoCo \rf{ver.mujoco} with a Panda arm and robosuite's
OSC\_POSE operational-space controller at \rf{ctrl.hz}\,Hz. In \emph{Lift} the robot grasps a
cube and raises it above a height threshold. In \emph{Can} (robosuite PickPlaceCan) it picks
a can from a source bin and places it in its target compartment of a second bin. Success is
robosuite's task predicate, evaluated on the current kinematics after every control step
(Section~\ref{sec:limitations}); an episode ends
at the first success or after \rf{proto.horizon} control steps. Initial states are
serialized MuJoCo snapshots (model XML, time, joint positions and velocities) that are
restored exactly for every policy, so no evaluation depends on a random reset.

\paragraph{Demonstrations.} We use the robomimic proficient-human (PH) low-dimensional
datasets for both tasks, in the version repackaged for robosuite~v1.5
(Hugging Face \code{robomimic/robomimic_datasets}, revision \rf{data.hf.revision}; file
SHA-256 prefixes \rf{data.lift.lowdim.sha} for Lift and \rf{data.can.lowdim.sha} for
Can)~\citep{mandlekar2021robomimic}. Each task has \rf{data.lift.demos} successful
demonstrations: Lift episodes last \rf{data.lift.len.min}--\rf{data.lift.len.max} control
steps (mean \rf{data.lift.len.mean}) and Can episodes
\rf{data.can.len.min}--\rf{data.can.len.max} (mean \rf{data.can.len.mean}). The data were
recorded through one teleoperation channel, so no claim about held-out operators is
possible.

\paragraph{Split.} Demonstrations are split by whole episode with robomimic's upstream
masks into \rf{data.lift.train} training and \rf{data.lift.valid} validation episodes, with
identical ID lists for both tasks. Normalization statistics are fit on the training episodes
only. Training windows never cross an episode boundary, and no observation after the
decision time is ever an input. The validation demonstrations give offline losses only: no
model, configuration or checkpoint is chosen with them. There is no demonstration test
split; the test is a set of held-out simulator initial states (Section~\ref{sec:protocol}).

\paragraph{Observations and privilege.} Every Level~1 policy receives the same
\rf{dim.obs}-dimensional vector per step: end-effector position and orientation
(quaternion), the two gripper finger joint positions, the object's position and orientation,
and the relative position of gripper and object. The object entries are exact simulator
poses, so every Level~1 result is a \emph{privileged-state} result and is never compared
causally with camera-based policies. For Can we compute the relative position from the
current step's poses; the dataset's stored relative-pose features lag by one step at the
first frame, so our Can input is not robomimic's default feature set and our Can numbers are
not comparable with numbers obtained on that set.

\paragraph{Actions.} The action is the controller's \rf{dim.act}-dimensional command: a
position delta, an axis-angle orientation delta and a gripper command, each scaled to
$[-1, 1]$ by the controller limits. The demonstrations' gripper commands are exactly
$\pm 1$. Every family predicts min--max normalized actions; predictions are mapped back with
the inverse transform and clipped to the action bounds by one shared decoder.

\paragraph{Platform.} Every reported rollout ran on one canonical evaluation platform: a
Linux aarch64 container (image \code{sha256:}\rf{plat.image}, Python \rf{ver.python}, OSMesa
rendering) on an Apple M1 Pro host, with policy inference on the CPU with one thread so that
rollouts are deterministic given their seeds. Physics is not bit-identical across platforms,
so native macOS rollouts were used for development only and never reported.
